\documentclass[aps,preprint]{revtex4}%
\usepackage{amsfonts}
\usepackage{amsmath}
\usepackage{amssymb}
\usepackage{graphicx}%
\setcounter{MaxMatrixCols}{30}
%TCIDATA{OutputFilter=latex2.dll}
%TCIDATA{Version=4.00.0.2321}
%TCIDATA{CSTFile=revtex4.cst}
%TCIDATA{Created=Friday, April 26, 2002 14:20:31}
%TCIDATA{LastRevised=Friday, April 26, 2002 16:15:36}
%TCIDATA{<META NAME="GraphicsSave" CONTENT="32">}
%TCIDATA{<META NAME="DocumentShell" CONTENT="Articles\SW\REVTeX 4">}
\newtheorem{theorem}{Theorem}
\newtheorem{acknowledgement}[theorem]{Acknowledgement}
\newtheorem{algorithm}[theorem]{Algorithm}
\newtheorem{axiom}[theorem]{Axiom}
\newtheorem{claim}[theorem]{Claim}
\newtheorem{conclusion}[theorem]{Conclusion}
\newtheorem{condition}[theorem]{Condition}
\newtheorem{conjecture}[theorem]{Conjecture}
\newtheorem{corollary}[theorem]{Corollary}
\newtheorem{criterion}[theorem]{Criterion}
\newtheorem{definition}[theorem]{Definition}
\newtheorem{example}[theorem]{Example}
\newtheorem{exercise}[theorem]{Exercise}
\newtheorem{lemma}[theorem]{Lemma}
\newtheorem{notation}[theorem]{Notation}
\newtheorem{problem}[theorem]{Problem}
\newtheorem{proposition}[theorem]{Proposition}
\newtheorem{remark}[theorem]{Remark}
\newtheorem{solution}[theorem]{Solution}
\newtheorem{summary}[theorem]{Summary}
\newenvironment{proof}[1][Proof]{\noindent\textbf{#1.} }{\ \rule{0.5em}{0.5em}}
\begin{document}
\preprint{HEP/123-qed}
\title[Short title for running header]{Long title that appears on the title page}
\author{First Author}
\affiliation{My Institution}
\author{Second Author}
\affiliation{My Institution}
\author{Third Author}
\affiliation{Other Institution}
\keywords{one two three}
\pacs{PACS number}

\begin{abstract}
Shell document for REV\TeX{} 4.

\end{abstract}
\volumeyear{year}
\volumenumber{number}
\issuenumber{number}
\eid{identifier}
\date[Date text]{date}
\received[Received text]{date}

\revised[Revised text]{date}

\accepted[Accepted text]{date}

\published[Published text]{date}

\startpage{101}
\endpage{102}
\maketitle
\tableofcontents


\section{Cones and Convex Sets}

A convex set is the convex closure of its extreme points in the relevant
topology i.e.%
\begin{equation}
S=\overline{Ext}\left\{  S\right\}
\end{equation}
where $x$ is an extreme point
\begin{align}
x &  =\left(  1-\alpha\right)  x_{1}+\alpha x_{2}\\
&  \Longrightarrow x_{2}=0\text{ and }x_{1}=x\\
i.e\quad\alpha &  =0
\end{align}
Then any element of the convex set can be expanded as
\begin{align}
x &  =\sum\alpha_{i}x_{i}\\
\sum\alpha_{i} &  =1;\quad\alpha_{i}\geq0\\
x_{i} &  \in Ext\left\{  S\right\}
\end{align}
The coefficients $\left\{  \alpha_{i}\right\}  $ are not unique. The Cone $C$
generated by $S$ is given by
\begin{equation}
C=\left\{  \alpha x|\forall\alpha\geq0,\quad x\in S\right\}
\end{equation}
Thus $y\in C$ if
\begin{align}
y &  =\alpha\sum\alpha_{i}x_{i}=\sum\alpha\alpha_{i}x_{i}=\sum\beta_{i}x_{i}\\
\sum\beta_{i} &  =\alpha\geq0,\quad\beta_{i}\geq0
\end{align}
Hence the extreme points of $S$ generate $C$. $S$ is called the base of the
cone $C$.


\end{document}